\subsection{Flare Candidate Identification}\label{sec:flare_detection}

We identified stellar flare candidates in the TESS photometry of
GJ~1132 using a semi-automated detection pipeline followed by visual
classification.  The pipeline ingests all available TESS 2-minute
cadence lightcurves from the Science Processing Operations Center
\citep[SPOC;][]{Jenkins2016} via the \texttt{lightkurve} package
\citep{Lightkurve2018}, covering Sectors~9, 10, 36, 63, 90, and~99.

\subsubsection{Preprocessing}

Each sector's lightcurve is first normalized to unit median flux using
the \texttt{lightkurve} \texttt{normalize()} method.  We then remove
low-frequency stellar variability by applying a Savitzky--Golay filter
\citep{Savitzky1964} with a window length of 101~cadences
($\approx$3.4~hr at 2-minute cadence).  The \texttt{lightkurve}
implementation iteratively sigma-clips outliers during the fit,
preventing flare signals from biasing the trend removal
\citep{Lightkurve2018}.  The resulting flattened lightcurve isolates
short-duration transient events against a unit baseline.

\subsubsection{Noise Estimation}

We estimate the photometric noise in each sector using the Median
Absolute Deviation (MAD), a robust measure of statistical dispersion
that is insensitive to outliers such as flares
\citep{Rousseeuw1993}.  For a sample of $N$ flux measurements
$\{F_i\}$, the MAD is defined as
\begin{equation}\label{eq:mad}
  \mathrm{MAD} = \mathrm{median}\bigl(|F_i - \widetilde{F}|\bigr),
\end{equation}
where $\widetilde{F} = \mathrm{median}(F_i)$.  The MAD is related to
the standard deviation of a Gaussian distribution by the consistency
constant $c = 1/\Phi^{-1}(3/4) \approx 1.4826$, where $\Phi^{-1}$ is
the quantile function of the standard normal distribution.  The
MAD-based robust sigma is therefore
\begin{equation}\label{eq:robust_sigma}
  \hat{\sigma} = 1.4826 \times \mathrm{MAD}.
\end{equation}
Unlike the sample standard deviation, $\hat{\sigma}$ has a breakdown
point of 50\%, meaning that up to half the data can be arbitrarily
corrupted without affecting the estimate \citep{Rousseeuw1993}.  This
property makes it well suited for flare detection, where bright
transient events would inflate a conventional $\sigma$ estimate and
suppress the detection of weaker flares.

\subsubsection{Candidate Detection}\label{sec:detection}

We flag all cadences in the flattened lightcurve whose flux exceeds
$\widetilde{F} + n_\sigma\,\hat{\sigma}$, where $n_\sigma$ is the
detection threshold.  We adopt $n_\sigma = 2.5$ to ensure sensitivity
to low-amplitude flares at the expense of a higher false-positive rate
that is subsequently filtered by visual inspection.  Individual
above-threshold cadences are grouped into contiguous clusters, allowing
a gap tolerance of up to 2 cadences ($\approx$4~min) to bridge
single-point dropouts within a flare event.  Clusters containing fewer
than $I_\mathrm{min} = 2$ consecutive above-threshold cadences are
discarded, as isolated single-point excursions are more likely due to
cosmic ray hits or instrumental artifacts than to stellar flares
\citep{Ilin2021}.

\subsubsection{Decay Tail Extension}\label{sec:decay_tail}

Stellar flares exhibit a characteristic fast-rise, exponential-decay
morphology \citep{Davenport2014}.  While the impulsive rise is
typically captured by the above-threshold cadences, the exponential
decay tail often extends well below the detection threshold.  To
capture the full flare energy, we extend each candidate window forward
in time from the last above-threshold cadence until the flux returns
to within $1\,\hat{\sigma}$ of the median for at least 3 consecutive
cadences ($\approx$6~min).  This criterion balances the need to
capture the decay phase against the risk of absorbing unrelated
variability.

\subsubsection{Candidate Characterization}

For each candidate, we record the following quantities:
\begin{itemize}
  \item \textbf{Peak time} ($t_\mathrm{peak}$): the TESS Barycentric
    Julian Date (BTJD) of maximum flux in the candidate window.
  \item \textbf{Peak significance} ($n_\sigma$): the peak flux
    expressed in units of $\hat{\sigma}$ above the sector median.
  \item \textbf{Duration}: the time interval from the first
    above-threshold cadence to the end of the decay tail, in minutes.
  \item \textbf{Equivalent duration} ($\mathrm{ED}$): the
    time-integrated relative flux excess,
    \begin{equation}\label{eq:equiv_dur}
      \mathrm{ED} = \int_{t_\mathrm{start}}^{t_\mathrm{stop}}
        \bigl(F(t)/F_\mathrm{quiescent} - 1\bigr)\,dt,
    \end{equation}
    with units of seconds \citep{Hunt-Walker2012}.  The equivalent
    duration is proportional to the flare energy and independent of
    distance, making it a natural unit for constructing flare frequency
    distributions \citep{Hawley2014}.  We evaluate the integral
    numerically using the trapezoidal rule on the normalized
    (unflattened) lightcurve, as the flattening process can distort
    the absolute flux scale of long-duration events.
\end{itemize}

Candidates are sorted by peak significance in descending order, so
that the most prominent events are presented first during visual
classification.

\subsubsection{Visual Classification}\label{sec:classification}

Each candidate is presented to the observer as a lightcurve panel
spanning $[-0.05, +0.10]$~days ($\approx$[$-72$, $+144$]~min) relative
to the peak time.  The normalized (unflattened) flux is plotted as a
black line, with the candidate interval highlighted using scatter
points.  The observer classifies each candidate into one of three
categories:
\begin{enumerate}
  \item \textbf{Flare}: the event displays morphology consistent with
    a stellar flare, typically a fast rise followed by an exponential
    decay.
  \item \textbf{Not a flare}: the event is attributed to instrumental
    artifacts, cosmic rays, or other non-astrophysical sources.
  \item \textbf{Uncertain}: the event is ambiguous and cannot be
    confidently classified.
\end{enumerate}

The classification is saved to a JSON file after each label, providing
crash-safe persistence.  The observer may navigate backward to
reclassify previous candidates and may quit the session at any point,
resuming later from the first unlabeled candidate.

\subsubsection{Pipeline Integration}

The labeled flare catalog is consumed by the downstream analysis
pipeline, which loads the JSON file and extracts the time intervals
of confirmed flares.  These intervals replace the previously
hard-coded flare parameters, ensuring that the flare frequency
distribution and energy analysis are derived from a systematic,
reproducible identification process.  The pipeline validates the
existence of the labeled catalog as a prerequisite before proceeding
with flare energy calculations and FFD construction.

\subsection{Joint Statistical Tension Between Kepler-Predicted and
TESS-Observed Flare Frequency Distributions}
\label{sec:joint_tension}

We quantify the discrepancy between the flare frequency distribution
(FFD) measured from TESS observations of GJ~1132 and the FFD
predicted by a Kepler-calibrated model evaluated at the stellar mass
and age of GJ~1132.  The FFD is parameterized as a power law in
log-space,
\begin{equation}
\log_{10} R = \alpha \log_{10} E + \beta,
\label{eq:ffd_powerlaw}
\end{equation}
where $R$ is the cumulative flare rate (day$^{-1}$) above energy $E$
(erg), $\alpha$ is the power-law slope, and $\beta$ is the intercept.
Both the TESS measurement and the Kepler prediction yield estimates of
the two-dimensional parameter vector $\boldsymbol{\theta} = (\alpha,
\beta)$.

\subsubsection{The TESS Measurement}

We fit Equation~\ref{eq:ffd_powerlaw} to the four TESS flares using
weighted least squares (\texttt{scipy.optimize.curve\_fit} with
\texttt{absolute\_sigma=True}), obtaining
\begin{equation}
\hat{\boldsymbol{\theta}}_{\mathrm{TESS}} =
  (\hat{\alpha}_{\mathrm{TESS}},\; \hat{\beta}_{\mathrm{TESS}}) =
  (-0.4498,\; 12.4864).
\end{equation}
The fit also returns the $2 \times 2$ covariance matrix
$\mathbf{C}_{\mathrm{TESS}}$, which encodes the uncertainties in
$\alpha$ and $\beta$ \emph{and their correlation}.  Because the four
flares span barely one order of magnitude in energy ($31.1 \lesssim
\log_{10} E \lesssim 32.2$), the slope and intercept are almost
perfectly anticorrelated: a slightly steeper slope can be compensated
by a higher intercept without moving the fit line away from the data.
The correlation coefficient is
\begin{equation}
r_{\alpha\beta} =
  \frac{C_{\alpha\beta}}{\sigma_\alpha \, \sigma_\beta} = -0.9999,
\end{equation}
where $\sigma_\alpha = 0.7205$ and $\sigma_\beta = 22.6824$ are the
marginal (one-dimensional) standard deviations read from the diagonal
of $\mathbf{C}_{\mathrm{TESS}}$.  This near-perfect anticorrelation
means the joint constraint occupies a razor-thin ellipse in $(\alpha,
\beta)$ space, even though the marginal uncertainties on each
parameter individually appear large.

\subsubsection{The Kepler Prediction}

The Kepler-calibrated model of \citet{Davenport2019} parameterizes the
FFD slope and intercept as linear functions of stellar age and mass:
\begin{align}
\alpha &= a_1 \log_{10} t + a_2 m + a_3,
  \label{eq:alpha_model} \\
\beta  &= b_1 \log_{10} t + b_2 m + b_3,
  \label{eq:beta_model}
\end{align}
where $t$ is the stellar age in Myr and $m$ is the stellar mass in
solar masses.  The six parameters $(a_1, a_2, a_3, b_1, b_2, b_3)$
were inferred via MCMC from a Kepler flare ensemble (see
Section~\ref{sec:kepler_ffd}), yielding posterior medians
$\boldsymbol{p}$ and a $6 \times 6$ covariance matrix
$\mathbf{C}_{6}$.

To predict $(\alpha, \beta)$ for GJ~1132, we evaluate
Equations~\ref{eq:alpha_model}--\ref{eq:beta_model} at $\log_{10} t =
\log_{10}(8000) \approx 3.903$ (i.e., 8~Gyr) and $m =
0.2\;M_\odot$:
\begin{equation}
\hat{\boldsymbol{\theta}}_{\mathrm{Kepler}} =
  (\hat{\alpha}_{\mathrm{Kep}},\; \hat{\beta}_{\mathrm{Kep}}) =
  (-1.0928,\; 34.4459).
\end{equation}
Because the model (Equations~\ref{eq:alpha_model}--\ref{eq:beta_model})
is linear in the six parameters, the uncertainty in the prediction is
obtained exactly via the Jacobian projection.  The $2 \times 6$
Jacobian matrix is
\begin{equation}
\mathbf{J} = \begin{pmatrix}
\log_{10} t & m & 1 & 0 & 0 & 0 \\
0 & 0 & 0 & \log_{10} t & m & 1
\end{pmatrix},
\label{eq:jacobian}
\end{equation}
and the projected $2 \times 2$ covariance of the Kepler prediction in
$(\alpha, \beta)$ space is
\begin{equation}
\mathbf{C}_{\mathrm{Kepler}} =
  \mathbf{J} \, \mathbf{C}_{6} \, \mathbf{J}^{\mathsf{T}}.
\label{eq:projected_cov}
\end{equation}
With 25{,}600 MCMC samples constraining six parameters,
$\mathbf{C}_{\mathrm{Kepler}}$ is very small---the Kepler prediction
is tight.

\subsubsection{Why Marginal Tensions are Misleading}

The simplest approach to quantifying the discrepancy is to compare the
predicted and observed values of each parameter separately:
\begin{equation}
n_{\sigma,\alpha} =
  \frac{|\hat{\alpha}_{\mathrm{TESS}} - \hat{\alpha}_{\mathrm{Kep}}|}
       {\sqrt{\sigma^2_{\alpha,\mathrm{Kep}} +
              \sigma^2_{\alpha,\mathrm{TESS}}}} = 0.89\;\sigma,
\qquad
n_{\sigma,\beta} =
  \frac{|\hat{\beta}_{\mathrm{TESS}} - \hat{\beta}_{\mathrm{Kep}}|}
       {\sqrt{\sigma^2_{\beta,\mathrm{Kep}} +
              \sigma^2_{\beta,\mathrm{TESS}}}} = 0.97\;\sigma.
\label{eq:marginal_tensions}
\end{equation}
These marginal tensions use only the diagonal elements of the
covariance matrices and therefore treat $\alpha$ and $\beta$ as if
they were independent.  When the correlation is weak, this is a
reasonable approximation.  However, with $r_{\alpha\beta} = -0.9999$,
the marginal standard deviations $\sigma_\alpha$ and $\sigma_\beta$
reflect the \emph{full extent} of the thin confidence ellipse
projected onto each axis.  These projections vastly overstate the
volume of parameter space consistent with the data.

Concretely, $\sigma_\alpha = 0.72$ and $\sigma_\beta = 22.7$ suggest
that the TESS fit is poorly constrained in both parameters.  But the
near-perfect anticorrelation means that knowing $\alpha$ determines
$\beta$ almost exactly (and vice versa).  The allowed region in
$(\alpha, \beta)$ space is not a $0.72 \times 22.7$ rectangle; it is
a sliver whose width perpendicular to its major axis is orders of
magnitude smaller.

\subsubsection{The Joint 2D Chi-Squared Test}

The correct approach accounts for the full $2 \times 2$ covariance
structure.  We define the residual vector
\begin{equation}
\boldsymbol{\Delta} =
  \hat{\boldsymbol{\theta}}_{\mathrm{TESS}} -
  \hat{\boldsymbol{\theta}}_{\mathrm{Kepler}} =
  (\hat{\alpha}_{\mathrm{TESS}} - \hat{\alpha}_{\mathrm{Kep}},\;
   \hat{\beta}_{\mathrm{TESS}} - \hat{\beta}_{\mathrm{Kep}}).
\label{eq:delta}
\end{equation}
Under the null hypothesis that both measurements estimate the same
underlying $(\alpha, \beta)$, and assuming the Kepler and TESS
measurements are independent (they use different instruments and
different stellar samples), the total covariance of the difference is
the sum of the individual covariances:
\begin{equation}
\mathbf{C}_{\mathrm{total}} =
  \mathbf{C}_{\mathrm{Kepler}} + \mathbf{C}_{\mathrm{TESS}}.
\label{eq:total_cov}
\end{equation}
Note that this sum preserves the off-diagonal terms.  Since
$\mathbf{C}_{\mathrm{Kepler}}$ is very small,
$\mathbf{C}_{\mathrm{total}} \approx \mathbf{C}_{\mathrm{TESS}}$, and
the strong anticorrelation in the TESS fit dominates.

The joint chi-squared statistic is then
\begin{equation}
\chi^2 = \boldsymbol{\Delta}^{\mathsf{T}} \,
  \mathbf{C}_{\mathrm{total}}^{-1} \, \boldsymbol{\Delta}.
\label{eq:chi2_joint}
\end{equation}
Under the null hypothesis, $\chi^2$ follows a chi-squared distribution
with 2 degrees of freedom (one for each component of
$\boldsymbol{\Delta}$).  This statistic measures the squared
Mahalanobis distance between the two estimates, properly accounting
for the shape and orientation of the confidence ellipse.

\subsubsection{Conversion to Equivalent Sigma}

To express the result as an equivalent Gaussian significance, we
compute the $p$-value from the $\chi^2$ distribution:
\begin{equation}
p = 1 - F_{\chi^2}(\chi^2;\, \nu = 2),
\label{eq:pvalue}
\end{equation}
where $F_{\chi^2}(\cdot;\, \nu)$ is the cumulative distribution
function of the chi-squared distribution with $\nu$ degrees of
freedom.  We then convert to an equivalent number of standard
deviations using the inverse of the standard normal survival function:
\begin{equation}
n_\sigma = \Phi^{-1}\!\left(1 - \frac{p}{2}\right),
\label{eq:nsigma}
\end{equation}
where $\Phi^{-1}$ is the quantile function (inverse CDF) of the
standard normal distribution.  The factor of 2 in the denominator
accounts for the two-tailed convention: $n_\sigma$ is the number of
standard deviations at which the two-tailed probability of a standard
normal equals $p$.

\subsubsection{Result}

Evaluating the above procedure yields
\begin{align}
\chi^2 &= 53.72 \quad (2\;\mathrm{DOF}), \\
p &< 10^{-4}, \\
n_\sigma &= 7.0.
\end{align}
The Kepler-predicted FFD for a $0.2\;M_\odot$ star at 8~Gyr is
discrepant with the TESS-observed FFD of GJ~1132 at $7.0\sigma$
significance.  GJ~1132 flares significantly less frequently than the
Kepler ensemble model predicts.

\subsubsection{Geometric Interpretation}

The geometry of the discrepancy is illustrated in the $(\alpha,
\beta)$ scatter plot (Figure~\ref{fig:fit_comparison}).  The Kepler
prediction occupies a tight cluster near $(\alpha, \beta) \approx
(-1.09, 34.4)$.  The TESS measurement, despite its large marginal
uncertainties, is confined to a thin diagonal band running from
upper-left to lower-right---the locus of lines that pass through the
observed flare data.  The Kepler prediction falls far from this band
in the direction perpendicular to it, which is why the joint tension
is so much larger than either marginal tension.

In the FFD plot itself (Figure~\ref{fig:ffd_comparison}), this
manifests as the Kepler prediction line (purple) lying well above the
grey $1\sigma$ confidence band of the TESS fit within the energy range
of the data.  The grey band is narrow despite the large marginal
uncertainties precisely because of the strong $\alpha$--$\beta$
anticorrelation: the band represents the family of lines consistent
with the data, and that family excludes the steeper, higher-intercept
Kepler prediction.
